\begin{examplebox}
	\textbf{\textcolor{CExample}{例1（基础）}}

	\textbf{\textcolor{CExample}{题目：}}
	已知函数$f(x)=x^2-2x$在区间$[0,3]$上，若不等式$a\ge f(x)$恒成立，求$a$的取值范围.

	\textbf{\textcolor{CExample}{解题步骤：}}
	\begin{description}[style=nextline, leftmargin=2.8em, labelsep=0.5em, itemsep=0.45em]
		\item[\textbf{第一步（求最大值）：}] 计算$f(x)$在$[0,3]$上的最大值。$f(x)=(x-1)^2-1$，在$[0,3]$上，$f(0)=0$，$f(1)=-1$，$f(3)=3$，故最大值为$3$。
			\[
				\max_{x\in[0,3]}f(x)=3.
			\]
			\par\textbf{理由：}二次函数开口向上，对称轴$x=1$在区间内，最大值在端点$x=3$处。

		\item[\textbf{第二步（套用结论一）：}] 由结论一，$a\ge f(x)$恒成立等价于$a\ge \max f(x)=3$。
			\[
				a\ge f(x)\ (\forall x\in[0,3]) \iff a\ge 3.
			\]
			\par\textbf{理由：}参数已分离，直接应用最值转化。

		\item[\textbf{第三步（写出范围）：}] 所以$a$的取值范围是$[3,+\infty)$。
			\[
				a\ge 3.
			\]
			\par\textbf{理由：}即满足条件的最小$a$为$3$。
	\end{description}
	\textbf{\textcolor{CExample}{关键结论：}}
	\textbf{$a\ge 3$}.

	\medskip
	\textbf{\textcolor{CExample}{例2（稍变形）}}

	\textbf{\textcolor{CExample}{题目：}}
	已知函数$f(x)=x+\frac{1}{x}$在区间$[1,3]$上，若不等式$a\le f(x)$恒成立，求$a$的取值范围.

	\textbf{\textcolor{CExample}{解题步骤：}}
	\begin{description}[style=nextline, leftmargin=2.8em, labelsep=0.5em, itemsep=0.45em]
		\item[\textbf{第一步（求最小值）：}] 由基本不等式，$x+\frac{1}{x}\ge2$，等号在$x=1$时成立。因为$1\in[1,3]$，所以最小值$2$。
			\[
				x+\frac{1}{x}\ge2\quad (x>0),
			\]
			等号在$x=1$时成立，故$\min_{x\in[1,3]}f(x)=2$。
			\par\textbf{理由：}基本不等式满足一正二定三相等，等点$1$在区间内。

		\item[\textbf{第二步（套用结论二）：}] 由结论二，$a\le f(x)$恒成立等价于$a\le \min f(x)=2$。
			\[
				a\le f(x)\ (\forall x\in[1,3]) \iff a\le 2.
			\]
			\par\textbf{理由：}参数分离后直接转化。

		\item[\textbf{第三步（写出范围）：}] $a\le 2$。
			\[
				a\le 2.
			\]
			\par\textbf{理由：}即$a$最大值为$2$。
	\end{description}
	\textbf{\textcolor{CExample}{关键结论：}}
	\textbf{$a\le 2$}.

	\medskip
	\textbf{\textcolor{CExample}{例3（综合应用）}}

	\textbf{\textcolor{CExample}{题目：}}
	已知不等式$2x^2+ax+1\ge0$对$x\in[1,2]$恒成立，求$a$的取值范围.

	\textbf{\textcolor{CExample}{解题步骤：}}
	\begin{description}[style=nextline, leftmargin=2.8em, labelsep=0.5em, itemsep=0.45em]
		\item[\textbf{第一步（分离参数）：}] 当$x\in[1,2]$时，$x>0$，原不等式等价于$a\ge -2x-\frac{1}{x}$。
			\[
				\begin{aligned}
					2x^2+ax+1\ge0 &\iff ax\ge -2x^2-1 \\
					&\iff a\ge -2x-\frac{1}{x}\quad (x>0).
			\end{aligned}
			\]
			\par\textbf{理由：}移项后除以$x>0$，不等号方向不变。

		\item[\textbf{第二步（求最值）：}] 令$g(x)=-2x-\frac{1}{x}$，它在$[1,2]$上递减，所以最大值在$x=1$处，$g(1)=-3$。
			\[
				\begin{aligned}
					g(x) &= -2x - \frac{1}{x}, \\
					\max_{x\in[1,2]}g(x) &= g(1), \\
					g(1) &= -3.
			\end{aligned}
			\]
			\par\textbf{理由：}由$y=-2x$和$y=-\frac{1}{x}$在$[1,2]$上都递减，故$g(x)$递减。

		\item[\textbf{第三步（套用结论一）：}] 由结论一，$a\ge g(x)$恒成立等价于$a\ge \max g(x)=-3$。
			\[
				a\ge g(x)\ (\forall x\in[1,2]) \iff a\ge -3.
			\]
			\par\textbf{理由：}此时参数$a$已分离，$g(x)$最值已知。
	\end{description}
	\textbf{\textcolor{CExample}{关键结论：}}
	\textbf{$a\ge -3$}.
\end{examplebox}
